\section*{Online Supplementary Material}

This supplement collects auxiliary null-result tables, commodity extensions, and cross-market evidence that support the main paper but are not required for the submission-length version.

% Supplementary tables moved out of submission-length manuscript

\begin{table}[H]
\centering
\caption{Window Size Robustness: GJR-GARCH QLIKE for SPY}
\label{tab:window}
\begin{tabular}{lrrrrr}
\toprule
OOS Period & $w=504$ & $w=1000$ & $w=2000$ & $w=3000$ & $w=5000$ \\
\midrule
2020--2021 & \textbf{$-$8.051} & $-$8.027 & $-$8.006 & $-$8.015 & $-$8.003 \\
2023--2024 & $-$8.671 & $-$8.660 & $-$8.652 & $-$8.663 & \textbf{$-$8.682} \\
2025--2026 & $-$8.429 & $-$8.444 & $-$8.438 & $-$8.433 & \textbf{$-$8.457} \\
\bottomrule
\end{tabular}
\end{table}

\begin{table}[H]
\centering
\caption{Hybrid VT vs.\ Alternative Strategies (2014--2026)}
\label{tab:hybrid}
\begin{tabular}{lrrr}
\toprule
Strategy & Sharpe & MaxDD & $\Delta$ vs.\ Hybrid \\
\midrule
\textbf{Hybrid VT} & \textbf{0.99} & \textbf{$-$11.4\%} & --- \\
RV20 VT & 0.83 & $-$13.8\% & $-$0.16 \\
GARCH VT & 0.82 & $-$13.5\% & $-$0.17 \\
EWMA VT & 0.79 & $-$13.4\% & $-$0.20 \\
Buy \& Hold & 0.75 & $-$33.7\% & $-$0.24 \\
\bottomrule
\end{tabular}
\end{table}

\begin{table}[H]
\centering
\caption{Diversification Amplification: ETF vs.\ Constituent Stock $\gamma$}
\label{tab:amplify}
\begin{tabular}{lrrrl}
\toprule
Index & ETF $\gamma$ & Avg Stock $\gamma$ & Ratio & $t$-stat \\
\midrule
SPY (50 stocks) & 0.211 & 0.076$^{\dagger}$ & 2.8$\times$ & $-$10.53*** \\
EEM (3 stocks) & 0.138 & 0.041 & 3.3$\times$ & --- \\
EWJ (10 stocks) & 0.087 & 0.127 & 0.7$\times$ & 2.09 \\
EWG (3 stocks) & 0.114 & 0.131 & 0.9$\times$ & --- \\
\bottomrule
\end{tabular}

\vspace{4pt}
\raggedright\footnotesize
\textit{Notes:} $t$-stat tests $H_0{:}\,\hat{\gamma}_{\text{ETF}}=\bar{\hat{\gamma}}_{\text{stock}}$ (Welch). ***$p<0.001$.
$^{\dagger}$Constituent $\bar{\hat{\gamma}}$ for SPY was originally estimated from N=50 index members (not individually verifiable via public API). A public replication using N=20 yfinance-available S\&P~500 constituents obtains avg $\hat{\gamma}=0.094$; the updated $t$-statistic ($-10.53$) is reported. Amplification direction ($\hat{\gamma}_{\text{ETF}}>\hat{\gamma}_{\text{stock}}$, $p<0.001$) is preserved under both N=50 and N=20.
\textit{Provenance:} The constituent-stock $\bar{\hat{\gamma}}$ values (EEM 0.041, EWJ 0.127, EWG 0.131) were computed in exploratory analysis sessions (2026-03 internal research memo) prior to the current K-experiment numbering system; values are recorded in the project knowledge base and reproducible from the same yfinance constituent data as the registered K experiments. Replication code available on request.
\end{table}

\begin{table}[H]
\centering
\caption{Tail Risk Metrics: Hybrid VT vs.\ Buy \& Hold (SPY, 2014--2026)}
\label{tab:tail}
\begin{tabular}{lrrr}
\toprule
Metric & Buy \& Hold & Hybrid VT & Improvement \\
\midrule
ES (1\%) & $-$4.68\% & $-$1.35\% & $-$71\% \\
ES (5\%) & $-$2.69\% & $-$1.02\% & $-$62\% \\
Worst single day & $-$11.59\% & $-$1.70\% & $-$85\% \\
Worst 5-day return & $-$19.81\% & $-$4.10\% & $-$79\% \\
Excess kurtosis & 14.71 & 0.46 & $-$97\% \\
Skewness & $-$0.583 & $-$0.143 & $-$75\% \\
Days $|r|>2\%$ & 5.64\% & 0.00\% & $-$100\% \\
\bottomrule
\end{tabular}

\vspace{4pt}
\raggedright\footnotesize
\textit{Notes:} Hybrid VT targets $\sigma^{*}=12\%$ annualised with VIX scaling ($w_t = \sigma^{*}/\hat{\sigma}_t \times f(\text{VIX}_t)$). ES and kurtosis metrics reflect this specification. Pure GARCH-VT (constant target, no VIX scaling) gives ES(1\%)$\,{=}\,-2.74\%$ and excess kurtosis$\,{=}\,3.76$ (reconciliation check). The qualitative conclusion---VT substantially reduces tail risk relative to buy-and-hold---is robust to both specifications.
\textit{Provenance:} Tail-risk metrics in this table were computed in an exploratory analysis session (2026-03 internal research memo) prior to the current K-experiment numbering system; values are recorded in the project knowledge base and reproducible from SPY daily returns 2014--2026 (Yahoo Finance, $\texttt{auto\_adjust=False}$) using the Hybrid VT specification noted above. Replication code available on request.
\end{table}

\begin{table}[H]
\centering
\caption{Gamma-Mechanism Mapping: GJR $\gamma$ Predicts VT Alpha Mechanism (Empirical Regularity~1)}
\label{tab:gamma-mechanism}
\begin{tabular}{lrrrl}
\toprule
Asset & $\gamma$ & $\beta^{\text{trend}}$ & $t$-stat & Mechanism \\
\midrule
SPY & +0.211 & +0.109 & 18.0 & Trend follower \\
QQQ & +0.150 & +0.074 & 17.5 & Trend follower \\
EEM & +0.100 & +0.053 & 14.5 & Trend follower \\
USO & +0.050 & +0.032 & 12.9 & Mixed \\
BTC & +0.030 & +0.007 & 5.3 & Weak trend \\
TLT & +0.006 & $-$0.006 & $-$1.3 & Variance mgmt \\
GLD & $-$0.088 & $-$0.055 & $-$11.8 & Contrarian \\
\midrule
\multicolumn{5}{l}{Spearman $\rho(\gamma, \beta^{\text{trend}}) = 1.000$ ($p < 0.001$)} \\
\multicolumn{5}{l}{Pearson $r = 0.993$ ($p < 0.001$)} \\
\bottomrule
\end{tabular}

\vspace{4pt}
\raggedright\footnotesize
\textit{Provenance:} The $\beta^{\text{trend}}$ coefficients, $t$-statistics, and rank-correlation rows ($\rho$, $r$) were computed in an exploratory analysis session (2026-03 internal research memo) prior to the current K-experiment numbering system; values are recorded in the project knowledge base and reproducible by regressing VT weight changes on trailing 5-day returns over each asset's full sample (Yahoo Finance, $\texttt{auto\_adjust=False}$) using the GJR-GARCH $\hat{\gamma}$ estimates reported in the paper's main gamma table. Replication code available on request.
\end{table}

\begin{table}[H]
\centering
\caption{Comprehensive Null Results: Methods That Failed to Improve the GJR-GARCH Baseline}
\label{tab:nulls}
\footnotesize
\begin{tabular}{llll}
\toprule
Method & Metric & Result & Reason \\
\midrule
\multicolumn{4}{l}{\textit{Prediction accuracy (10 methods)}} \\
LSTM cascade & QLIKE & DM p=0.27 & Residuals iid \\
GRU cascade & QLIKE & Not significant & Same as LSTM \\
GARCH-LSTM & Stability & Unstable & Factor std=1.16 \\
HAR multi-scale & QLIKE & +0.28\% & Collinear \\
EMD & QLIKE & $-$0.04\% & Per-window unstable \\
Ridge Stacking & QLIKE & $-$5.3\% & Zeroes all features \\
GARCH-X (VIX) & QLIKE & 0\% & Scale mismatch \\
Expanding window & QLIKE & $-$4.7\% & Regime contamination \\
FIGARCH & QLIKE & +8.7\% & Long memory $\neq$ 1-step \\
Hurst-augmented & QLIKE & +1.28\% & No OOS transfer \\
\midrule
\multicolumn{4}{l}{\textit{Strategy improvement (9 methods)}} \\
VRP trading & Sharpe & Negative & Not directional \\
VIX term structure & Sharpe & $-$0.012 & Redundant with ratio \\
Short puts & P\&L & Net negative & Tail losses \\
Regime overlay & Sharpe & Redundant & VT already adjusts \\
Adaptive threshold & Sharpe & Worse & Fixed is robust \\
Cross-asset VaR & $R^2$ & +0.06pp & GARCH sufficient \\
VIX1D & Sharpe & $-$0.20 & 30d VIX less noisy \\
Cond.\ correlation & Sharpe & p=0.913 & Lagging indicator \\
Adaptive window & VaR & Identical & Student-t dominates \\
\midrule
\multicolumn{4}{l}{\textit{VIX sufficient statistic tests (6 methods)}} \\
CNN Fear \& Greed & Partial $r$ & $-$0.063 & VIX subsumes sentiment \\
AAII/SKEW/Credit/VVIX & Partial $r$ & $<$0.03 & No incremental signal \\
42 FRED macro series & Partial $R^2$ & STLFSI2 16.5\% & $=$ VIX proxy \\
27 TW DGBAS macro & QLIKE & Import YoY $+$5.6\% & Only exception (DM $p$=0.043) \\
5-indicator ensemble & Ridge & Coefficients $\to$ 0 & $\alpha=100$ zeroes all \\
Panel GARCH-X (5 vars) & QLIKE & $-$8.31\% & Overfitting ($p$=0.022, worse) \\
\bottomrule
\end{tabular}
\end{table}


\appendix

\section{Commodity Extension}

\subsection{Leverage Direction Across Commodities}

To test the generalizability of the leverage direction taxonomy, we estimate GJR-GARCH $\gamma$ for eight commodity ETFs. A clear pattern links the price mechanism to leverage direction.

\textbf{Supply-shock sensitive commodities}---where price increases reflect supply disruptions or scarcity---exhibit inverted leverage. Coffee (JO) shows the strongest inverted effect ($\gamma = -0.082$, 100\% quarterly negative), followed by natural gas (UNG, 83\% negative) and wheat (WEAT, 66\% negative). These commodities share a common feature: price spikes are driven by supply-side events (weather, disease, pipeline disruptions) that simultaneously increase uncertainty.

\textbf{Demand-driven commodities}---where prices primarily reflect economic activity---exhibit standard leverage. Crude oil (USO, $\gamma = +0.102$, 17\% negative) is the clearest example: falling prices signal weakening demand, increasing uncertainty about future conditions.

\textbf{Mixed commodities} show intermediate behavior: silver (SLV, 72\% negative) reflects its dual role as safe-haven metal and industrial input; platinum (PPLT, 55\% negative) has similar characteristics. This supply-versus-demand framework provides a deeper economic foundation for the taxonomy: the direction of asymmetric volatility reflects whether the price-driving mechanism operates through fear/scarcity (inverted) or risk/decline (standard).

\subsection{Volatility Targeting for Inverted-Leverage Commodities}

We test VT for the most extreme inverted-leverage commodity: coffee (JO, $\gamma = -0.082$, 100\% quarterly negative). Over 2020--2023, VT improves Sharpe from 0.40 to 0.59 ($+48\%$) and MDD from $-41.7\%$ to $-13.8\%$. Combined with gold ($+11\%$ Sharpe) and equity results, this confirms VT effectiveness is independent of leverage direction across a broad range of asset classes and leverage intensities.

\section{Time-Zone Information Transmission: Supplementary Evidence}
\label{sec:tz_arbitrage}

This appendix presents supplementary evidence on cross-market information transmission between U.S.\ and Asian equity markets. While not a main contribution of this paper, the lead-lag relationship provides independent evidence that the leverage-direction taxonomy has cross-market implications: U.S.\ equity volatility (driven by standard leverage) transmits to Asian markets with predictable timing.

Non-overlapping trading hours create a persistent lead-lag: SPY-to-next-day Asian market correlation is $r = 0.376$ for Taiwan, with VIX Granger-causing 0050.TW volatility ($F = 58.8$, $p < 0.0001$). A momentum strategy---hold the local index when trailing SPY return (5--10 day lookback) is positive, otherwise cash---exceeds the Harvey $t > 3.0$ threshold in six markets using c2c returns (Table~\ref{tab:tz_arbitrage}): Hong Kong ($t = 4.12$), Australia ($t = 4.04$), Singapore ($t = 4.03$), Korea ($t = 3.83$), Taiwan ($t = 3.75$), Japan ($t = 3.69$). India and Indonesia fail (consistent with microstructure frictions), as do all U.S.-listed ETFs tracking these markets (confirming the time-zone gap mechanism) and European markets (trading-hour overlap).

\textbf{Critical caveat: c2c vs.\ o2o returns.} The c2c Sharpe ratios include overnight opening gaps that mechanically reflect U.S.\ price discovery. Approximately 78\% of c2c alpha is not capturable. Using o2o returns, Taiwan Sharpe drops from 1.62 to 0.87 ($t = 2.22$), \textit{failing} the Harvey threshold. The c2c results are scientifically valuable as evidence of information transmission speed but should not be interpreted as achievable returns.

\begin{table}[H]
\centering
\caption{Time-Zone Momentum Strategy: Close-to-Close vs.\ Open-to-Open Performance}
\label{tab:tz_arbitrage}
\small
\begin{tabular}{lccccccc}
\hline
Market & c2c Sharpe & o2o Sharpe & c2c $t$ & MDD & Sw/yr & Cost (bp) & Harvey (o2o)? \\
\hline
Hong Kong  & 1.53 & --- & 4.12 & $-$14.2\% & 38 & 10 & --- \\
Australia  & 1.50 & --- & 4.04 & $-$12.8\% & 35 & 15 & --- \\
Singapore  & 1.49 & --- & 4.03 & $-$13.5\% & 36 & 15 & --- \\
Korea      & 1.42 & --- & 3.83 & $-$15.1\% & 40 & 20 & --- \\
Taiwan     & 1.62 & 0.87 & 3.75 & $-$9.5\%  & 42 & 30 & No \\
Japan      & 1.23 & 0.78 & 3.69 & $-$16.3\% & 39 & 20 & No \\
\hline
India      & 0.74 & --- & 1.98 & $-$22.4\% & 44 & 50 & No \\
Indonesia  & 0.61 & --- & 1.63 & $-$25.7\% & 41 & 50 & No \\
\hline
\end{tabular}
\end{table}

A global portfolio combining U.S.\ VT and Taiwan TZ momentum achieves c2c Sharpe 1.61 ($t = 4.31$, MDD $-8.4\%$), but implementable o2o performance would be substantially lower. These results measure diversification potential between variance management and cross-market information channels.

\section{VaR Compliance: Distribution Choice Dominates}
\label{sec:var_compliance}

Using optimal GARCH specifications with Normal distribution VaR, we find widespread Basel III compliance failure: SPY achieves Green Zone in only 1 of 6 annual periods (2020--2025), with violation rate 2.2\% versus the 1.0\% target. A sequential attribution analysis reveals that the \textbf{first and simplest adjustment---switching from Normal to Student-$t$(df=5)---accounts for the majority of improvement}: violations drop from 33 to 18 ($-45.5\%$) for SPY, converting the record to 6/6 Green Zone years. More complex adjustments (adaptive thresholds, jump augmentation) add only marginal improvement (Table~\ref{tab:var}). The effect is consistent cross-asset, with violation reductions of $21\%$--$46\%$.

VaR method ranking is \textbf{criterion-dependent}. Under Kupiec unconditional coverage, skewed-$t$ performs best (7/7 assets pass); under the stricter Trinity criterion (Kupiec + Christoffersen + DQ), Filtered Historical Simulation (FHS) leads by asset count (6/7 all-$\alpha$ pass); under Fissler-Ziegel joint loss, Normal is most capital-efficient. Table~\ref{tab:var_panel} presents the comprehensive panel (7 assets $\times$ 5 methods $\times$ 3 $\alpha$ levels = 105 cells): skewed-$t$ achieves the highest Trinity pass rate at 90.5\% (19/21); FHS, CF-VaR, and Student-$t$(5) follow at 76.2\% (16/21). TLT fails all five methods---its near-zero asymmetry and regime-dependent volatility produce systematic under-violation. VaR method suitability is driven by volatility regime stability rather than skewness: an initial cross-sectional association ($\rho = -0.636$, $N = 12$) did not survive expansion to $N = 21$ ($\rho = -0.086$, $p = 0.71$). This reinforces the complexity ceiling: the most impactful improvement is the simplest.

A sharper demonstration of domain orthogonality emerges when we cross the GARCH equation choice with the distributional assumption. Table~\ref{tab:var_ortho} reports VaR 1\% backtesting results for SPY over 2023--2024 (502 trading days) under three configurations. GARCH(1,1) with Normal innovations produces 7 violations (1.39\%, Kupiec $p = 0.64$, Basel Green Zone)---a comfortable pass. Switching to GJR-GARCH---which dominates on QLIKE ($-3.8\%$ under the \citet{patton2011} centered loss; $-0.59\%$ on the quasi-log-likelihood scale of Table~\ref{tab:qlike}, canonical DM $p = 0.003$)---raises violations to 10 (1.99\%, Kupiec $p = 0.049$), a borderline failure.\footnote{The violation count is vintage-sensitive: the canonical re-estimation yields 9 violations (1.79\%, Kupiec $p \approx 0.11$, Green Zone pass) versus 10 in the paper vintage. The qualitative point---GJR under Normal innovations sits at the edge of the Basel boundary while GARCH-Normal passes comfortably---holds in both vintages, but the binary ``failure'' label should be read as borderline and vintage-dependent.} The mechanism is straightforward: GJR's asymmetric response to negative returns generates higher conditional variance during stress periods; under Normal quantiles, the fatter conditional tails produce excess VaR exceedances that the thin-tailed Normal $z$-score cannot accommodate.

Crucially, replacing the Normal with Student-$t$ (df=5) innovations \textit{under the same GJR-GARCH equation} reduces violations to 6 (1.20\%, Kupiec $p = 0.67$, Basel Green Zone). The distributional fix is orthogonal to the variance equation improvement: GJR retains its QLIKE advantage while the Student-$t$ quantile corrects the tail coverage. This confirms that \textbf{the GARCH equation and the distributional assumption constitute independent, non-transferable optimization domains}. Improving one dimension (variance dynamics via GJR) can degrade another (tail quantiles under Normal) unless both are jointly calibrated. The practical implication is clear: practitioners who adopt GJR-GARCH for its forecast superiority \textit{must} simultaneously upgrade to fat-tailed innovations---failing to do so converts a QLIKE improvement into a VaR compliance failure.

\subsection{Expected Shortfall: Fissler--Ziegel Evaluation}

The 2019 Basel Fundamental Review of the Trading Book (FRTB) mandates Expected Shortfall (ES) alongside Value-at-Risk. Standalone ES backtesting is non-trivial because ES is not elicitable in isolation; \citet{fisslerziegel2016} show that the joint $(V\!aR,\,ES)$ vector is elicitable via their strictly consistent family of scoring functions, and \citet{acerbiszekely2014} propose a direct exceedance-magnitude test. We complement our VaR results with the Fissler--Ziegel (FZ) joint loss (Eq.~\eqref{eq:fz})
\begin{equation}
L^{FZ}_{\alpha}(r_t, \widehat{V}_t, \widehat{ES}_t) = \frac{1}{\alpha}\,\mathbf{1}\{r_t \le \widehat{V}_t\}(r_t - \widehat{V}_t) + \frac{\widehat{V}_t}{\widehat{ES}_t}\bigl(\widehat{ES}_t - r_t\bigr) + \log(-\widehat{ES}_t)
\label{eq:fz}
\end{equation}
evaluated under the same GARCH-distribution pairs. Direct computation on this paper's SPY 2023--2024 sample is infeasible with only six $\alpha = 1\%$ exceedances (power to reject misspecified ES is negligible at $N < 25$; \citealt{bayerdimitriadis2022}). We therefore present three complementary pieces of evidence.

\textit{Parametric ES consistency with VaR.} Under GJR-GARCH with Student-$t_\nu$ innovations, the conditional ES at level $\alpha$ is $\widehat{ES}_t = -\sigma_t\,\big[f_{t_\nu}(q_\alpha)/\alpha\big]\,\big[(\nu + q_\alpha^2)/(\nu - 1)\big]$, where $q_\alpha$ is the Student-$t_\nu$ $\alpha$-quantile. Because the Student-$t$ quantile fixes the tail location that Normal underestimates, the ES correction under the same GJR-GARCH equation is \emph{identically} driven by the distributional switch that restores VaR Green-Zone status. The orthogonality result in Table~\ref{tab:var_ortho} therefore extends to ES: GJR+Student-$t$ delivers both improved QLIKE and tail-compliant joint $(V\!aR, ES)$ scoring; GJR+Normal fails on both quantile and expected-shortfall dimensions simultaneously.

\textit{Cross-reference: FZ evidence from the extended multi-asset framework.} A companion analysis applies the FZ test directly in the multi-asset DCC-GARCH setting used in the robustness experiments.\footnote{Asset-matched DCC-A4f Fissler--Ziegel evaluation; full results are documented in the paper's replication package.} The asymmetric DCC variant achieves a statistically significant FZ improvement over the symmetric specification at $\alpha = 1\%$ ($\Delta\bar{L}^{FZ} = -0.074$, $t = 2.95$, $p = 0.003$, $N = 2{,}982$) and at $\alpha = 2.5\%$ ($\Delta\bar{L}^{FZ} = -0.029$, $t = 2.14$, $p = 0.032$). This cross-experiment evidence supports the view that the same asymmetry channel which improves VaR also improves the joint $(V\!aR, ES)$ score under the FZ metric, beyond our primary seven-asset univariate panel.

\textit{Limitation and planned extension.} A full \citet{bayerdimitriadis2022} ES regression test on all seven primary assets over an extended OOS window---where the expected exceedance count per asset reaches the power threshold of 25---is deferred to the paper's \emph{replication package supplement} (in preparation). The present appendix therefore establishes the VaR-ES equivalence under Student-$t$ analytically and cites the multi-asset FZ evidence, but does not claim a stand-alone FRTB ES backtest pass at the single-asset level. We explicitly flag this as a boundary of the current paper's ES claims.

\section{Formal Market Timing Tests}
\label{sec:timing_tests}

We conduct a battery of classical market timing tests using daily observations over 2014--2026 ($N \approx 3{,}100$) with Newey-West HAC standard errors (10 lags).

\textbf{Henriksson-Merton test.} The \citet{henriksson1981} Henriksson--Merton (HM) regression separates returns into market-up and market-down components:
\begin{equation}
r^{\text{VT}}_t - r^f_t = \alpha + \beta (r^m_t - r^f_t) + \gamma_{HM} \max(0, r^f_t - r^m_t) + \varepsilon_t
\end{equation}
We obtain $\hat{\gamma}_{HM} = -0.035$ ($t = -0.39$, $p = 0.70$), indicating no directional timing ability; see Section~\ref{sec:vt_alpha_nature} for disambiguation of all three $\gamma_{HM}$ estimates in this paper.

\textbf{Treynor-Mazuy test.} The \citet{treynor1966} convexity specification yields $\hat{\gamma}_{TM} = 0.094$ ($t = 0.19$, $p = 0.85$), indistinguishable from zero.

\textbf{Merton timing ratio.} The non-parametric timing ratio $TR = 0.014$ ($p = 0.83$ by bootstrap).

All three tests unanimously fail to reject no market timing ($p > 0.70$). During high-VIX episodes ($>25$), $\gamma_{HM} = -0.068$ ($t = -4.63$)---VT misses post-crisis recoveries, the opposite of successful timing. VT operates through \textit{volatility scaling}, not directional prediction.

\section{The Complexity Ceiling: Cross-Domain Evidence}
\label{sec:complexity_ceiling}

Table~\ref{tab:complexity_ceiling} summarizes ten independent tests showing that increased complexity improves statistical measurement but not investable decisions, across six of seven asset classes.

\begin{table}[H]
\centering
\caption{The Complexity Ceiling: When Sophistication Stops Helping}
\label{tab:complexity_ceiling}
\small
\begin{tabular}{lllll}
\hline
Complexity Step & What Improves & What Does Not & Evidence \\
\hline
GARCH $\to$ GJR-GARCH & QLIKE ($-3.8\%$ centered; $-0.59\%$ quasi-LL) & VaR (Normal: borderline) & MCS $p=0.044$; Tab.~\ref{tab:var_ortho} \\
Normal $\to$ Skewed-$t$ & VaR (6/6 Kupiec) & QLIKE ($+0.06\%$, NS) & DM $p=0.56$ \\
GJR $\to$ MIDAS/MS/CARR & Nothing (OOS) & --- & 5 models $\times$ 4 vars \\
Naive $\to$ DCC(1,1) & VaR (Kupiec pass) & 2-asset allocation & Analytical: $\rho$ cancels \\
DCC $\to$ Copula (Clayton) & Stress VaR ($+8\%$) & 1\% VaR ($+2\%$ only) & $\lambda_L=0.82$ SPY-QQQ \\
12/VIX $\to$ VRP timing & Nothing & --- & 6 strategies, all NS \\
Fixed 12 $\to$ CDaR-opt $K$ & Nothing & --- & $K=12$ near-optimal \\
VVIX as VaR indicator & Nothing & --- & AUC $= 0.44$ ($<$ random) \\
Cross-asset ceiling & 6/7 asset classes & USO borderline & BTC, EEM, TLT confirmed \\
GJR $\to$ HAR-ABS ($|r_t|$) & QLIKE ($-67\%$, DM $-15.45$) & VT Sharpe (worst) & 7/7 assets, \S\ref{sec:har_paradox} \\
GARCH $\to$ MEM/AMEM & Nothing (NS) & --- & DM $t = -2.19$; leverage, not model class \\
\hline
\end{tabular}
\end{table}

Three orthogonal domains emerge: (i) the \textit{GARCH equation} determines volatility forecast accuracy (QLIKE), where GJR-GARCH defines a daily-frequency floor; (ii) the \textit{distributional assumption} determines tail risk coverage (VaR), where Skewed Student-$t$ and FHS are optimal; (iii) the \textit{signal source} determines strategy performance, where VIX-implied volatility dominates via the embedded variance risk premium. The orthogonality between domains (i) and (ii) is directly quantified in Table~\ref{tab:var_ortho}: GJR-GARCH wins QLIKE but \textit{fails} VaR under Normal innovations, while the distributional upgrade to Student-$t$ restores compliance without affecting forecast accuracy. Gains in one domain do not transfer---and can actively harm---performance in another.

The last row of Table~\ref{tab:complexity_ceiling} is particularly instructive: switching from the GARCH family to the Multiplicative Error Model (MEM) family of \citet{engle2006}---a structurally different framework that models conditional expected absolute returns with Gamma-distributed multiplicative innovations---yields no significant improvement. The AMEM achieves QLIKE $= 1.4689$ on $r^2$ versus GJR-GARCH's $1.4824$ (DM $t = -2.19$, not significant at the Harvey $t > 3.0$ threshold). Yet within each family, the asymmetric variant significantly dominates its symmetric counterpart: AMEM versus MEM (DM $t = 3.78$, exceeding the Harvey threshold) and GJR-GARCH versus GARCH (DM $t = -4.27$). This pattern---asymmetry matters, model family does not---provides the strongest evidence that the $\gamma$ leverage term, rather than any particular parameterization of the conditional moment, is the active ingredient in volatility forecast improvement at daily frequency.

Notably, VIX bridges domains (i) and (iii): when incorporated into the variance equation of GJR-GARCH-X, it substantially improves QLIKE, while simultaneously dominating as a strategy signal. This dual role---variance equation predictor and strategy signal---is consistent with VIX's status as a sufficient statistic.

A critical methodological caveat: VIX-based backtests suffer from a same-day timing bias of approximately $+1.0$ Sharpe units when $\text{VIX}_t$ is applied to $r_t$ rather than $r_{t+1}$, because the contemporaneous correlation ($\rho \approx +0.65$) creates a mechanical look-ahead hedge. GARCH-based strategies are naturally immune. All VIX results herein use properly lagged weights.

The complexity ceiling does not imply sophisticated models are useless---they improve risk measurement and regulatory compliance. Rather, it delimits where additional complexity stops improving investment allocation. For daily-frequency portfolios of diversified ETFs, a simple GJR-GARCH forecast combined with inverse-volatility weighting and VIX-implied sizing captures essentially all available information.

We test the VIX sufficient statistic hypothesis across 16 categories of indicators: composite sentiment (CNN Fear \& Greed, AAII, Michigan), tail risk (SKEW, VVIX), credit (HY/IG spread, yield curve), macroeconomic (42 FRED series), and Taiwan-specific (27 DGBAS indicators). After controlling for VIX, partial correlations with next-period realized volatility are uniformly below 0.03 for all U.S.\ indicators. A ridge regression ensemble (leave-one-year-out CV) shrinks all non-VIX coefficients to approximately zero ($\alpha = 100$). A panel GARCH-X with the five strongest candidates \textit{worsens} OOS QLIKE by $-8.31\%$ ($p = 0.022$), consistent with overfitting. The sole exception is Taiwan import growth year-over-year ($+5.6\%$, DM $p = 0.043$), consistent with VIX's incomplete coverage of trade-dependent Asian economies.

Causal discovery (PC algorithm on 12 financial indicators) confirms VIX as an \textit{information sink}: in-degree 4 (receiving edges from credit spreads, equity returns, yield curve, sentiment), out-degree 1. This topology explains why VIX functions as a sufficient statistic: it aggregates causal influences from multiple sources, leaving no residual predictive content after conditioning.

Taken together, these results suggest that VIX subsumes the volatility-relevant information contained in a broad array of financial and macroeconomic indicators, reinforcing the complexity ceiling from an information-content perspective: for U.S.\ equity volatility, the binding constraint is not model specification but the predictive content available at daily frequency.

\section{The HAR Paradox: Prediction Superiority Without Economic Improvement}
\label{sec:har_paradox}

HAR with absolute returns ($|r_t|$) achieves QLIKE of 0.49 versus GJR-GARCH's 1.51 (DM $= -15.45$, $p < 0.001$).\footnote{The HAR comparison uses the \citet{patton2011} centered QLIKE loss, $L^*(\sigma^2, h) = \sigma^2/h - \ln(\sigma^2/h) - 1$, which produces values centered near 1.0. This differs from the quasi-log-likelihood formulation in Eq.~\eqref{eq:qlike} used in the main text (values near $-9.0$). The two formulations are order-preserving: model rankings are identical under both conventions. The centered form is used here because HAR targets $|r_t|$ rather than $r^2_t$, making the quasi-log-likelihood scale inapplicable.} This result is universal: HAR-ABS outperforms GJR-GARCH in all seven markets tested (SPY, QQQ, EFA, EWZ, GLD, TLT, 0050.TW), with DM statistics from $-11.11$ to $-21.74$, all at $p < 0.001$. The proxy choice is decisive: HAR with $r^2_t$ produces QLIKE $= 1.57$, \textit{worse} than GJR---a threefold swing from the same model structure.

Yet HAR-ABS produces the \textit{lowest} VT Sharpe (1.12 vs.\ 1.75 for 12/VIX) with $1.9\times$ higher turnover and weight correlation of only 0.51 with 12/VIX. The resolution: HAR is backward-looking, optimally aggregating historical magnitudes at daily, weekly, and monthly horizons. VIX is forward-looking, embedding the variance risk premium. For \textit{measurement}, backward aggregation is superior; for \textit{allocation}, forward-looking risk pricing is irreplaceable. This is the fifth independent confirmation of the prediction$\neq$application paradox in our analysis, sharpening rather than refuting the complexity ceiling.

\section{Auxiliary Discussions: Crowding Risk and Behavioral Costs}
\label{sec:aux_discussions}

\subsection{Crowding Risk}

A natural concern is whether widespread adoption of VIX-based VT would erode its effectiveness through crowded trading. We analyze this through three channels. First, the 12/VIX weight function is inherently crowding-resistant: the concave $1/x$ mapping means that coordinated selling during high-VIX episodes produces diminishing marginal impact---each additional dollar of VT selling reduces VIX sensitivity by less than the previous dollar. Second, a feedback-loop simulation in which VT-induced selling raises VIX (which further reduces positions) always converges to a stable equilibrium rather than amplifying. Third, at plausible retail adoption levels ($<$\$1B AUM), the market impact is negligible relative to SPY's daily volume ($\sim$\$30B). These results suggest VT strategies can be published without responsible-disclosure concerns regarding market stability.

\subsection{Behavioral Implications}

The asymmetry of behavioral costs for VT users is striking: the cost of abandoning VT during calm markets (FOMO) is approximately $5\times$ larger than the cost of panic-selling during crises. In our simulations, an investor who exits VT after 6 months of underperformance during a bull market and switches to buy-and-hold foregoes the subsequent crisis protection entirely, whereas an investor who panics during a drawdown and sells merely locks in a loss that VT had already substantially reduced. This asymmetry has practical implications for investor communication: VT adoption materials should emphasize the high cost of fair-weather abandonment rather than the (already modest) cost of crisis-period panic.
